\documentclass{beamer}
\usepackage{listings}
\newcommand\blank[1]{\rule[-.2ex]{#1}{.4pt}}
\usepackage{tikz}
\usetikzlibrary{shapes.callouts, positioning, arrows.meta,calc,backgrounds}
\usepackage[dvipsnames]{xcolor}

% Use plain theme with no decorations
\usetheme{default}
\setbeamertemplate{navigation symbols}{}
\setbeamertemplate{footline}{}
\setbeamertemplate{headline}{}

% Define Java code style
\lstdefinestyle{JavaStyle}{
    language=Java,
    basicstyle=\ttfamily\small,
    keywordstyle=\color{blue}\bfseries,
    commentstyle=\color{green!60!black},
    stringstyle=\color{Maroon},
    numbers=left,
    numberstyle=\tiny\color{gray},
    stepnumber=1,
    numbersep=10pt,
    breaklines=true,
    frame=single,
    tabsize=2,
    showstringspaces=false,
    firstnumber=auto
}

\title{Understanding Java Code Execution}
\subtitle{Step-by-Step Walkthrough}
\author{Programming Tutorial}
\date{}

\begin{document}

\begin{frame}
\titlepage
\end{frame}

\begin{frame}[fragile]{Code Execution Flow}
    \begin{center}
    \begin{tikzpicture}[remember picture, overlay]
        % Code listing (named node so we can reference its bounds)
        \node[anchor=center] (code) at (current page.center) {
            \hspace{1cm}
            \begin{minipage}{\textwidth}
    \begin{lstlisting}[style=JavaStyle,firstnumber=1]
    
public class Main {
    public static void main(String[] args) {
        System.out.println("Hello World: 1");
    \end{lstlisting}
            \end{minipage}
        };
        
        % === Spotlight the current line: draw a semi-opaque overlay with a hole ===
        \begin{scope}[on background layer]
            % Compute a window rectangle around the current line
            \path let \p1=(code.north west), \p2=(code.north east) in
                coordinate (winNW) at (\x1+0.9em, \y1-2.6em)  % top-left of the window
                coordinate (winSE) at (\x2-0.9em, \y1-1.5em); % bottom-right of the window
            
            % Wash out (pseudo-blur) everything except the window using even-odd rule
            \fill[black,opacity=0.5,even odd rule]
                (code.south west) rectangle (code.north east)
                (winNW) rectangle (winSE);
            
            % Draw the red box around the visible line
            \draw[red,very thick,rounded corners=2pt] (winNW) rectangle (winSE);
        \end{scope}
        
        % Explanation box below the code
        \node[anchor=north,fill=blue!10,draw=blue!50,thick,rounded corners,
            text width=0.8\textwidth,align=left,inner sep=10pt] 
            at ($(current page.north)+(0,-1.5cm)$) {
            \textbf{Line 1:} Define the Main class
        };
    \end{tikzpicture}
    \end{center}
    \end{frame}

    \begin{frame}[fragile]{Code Execution Flow}
    \begin{center}
    \begin{tikzpicture}[remember picture, overlay]
        % Code listing (named node so we can reference its bounds)
        \node[anchor=center] (code) at (current page.center) {
            \hspace{1cm}
            \begin{minipage}{\textwidth}
    \begin{lstlisting}[style=JavaStyle,firstnumber=1]
    
public class Main {
    public static void main(String[] args) {
        System.out.println("Hello World: 1");
        System.out.println("Hello World: 2");
    \end{lstlisting}
            \end{minipage}
        };
        
        % === Spotlight the current line: draw a semi-opaque overlay with a hole ===
        \begin{scope}[on background layer]
            % Compute a window rectangle around the current line
            \path let \p1=(code.north west), \p2=(code.north east) in
                coordinate (winNW) at (\x1+0.9em, \y1-3.7em)  % top-left of the window
                coordinate (winSE) at (\x2-0.9em, \y1-2.6em); % bottom-right of the window
            
            % Wash out (pseudo-blur) everything except the window using even-odd rule
            \fill[black,opacity=0.5,even odd rule]
                (code.south west) rectangle (code.north east)
                (winNW) rectangle (winSE);
            
            % Draw the red box around the visible line
            \draw[red,very thick,rounded corners=2pt] (winNW) rectangle (winSE);
        \end{scope}
        
        % Explanation box below the code
        \node[anchor=north,fill=blue!10,draw=blue!50,thick,rounded corners,
            text width=0.8\textwidth,align=left,inner sep=10pt] 
            at ($(current page.north)+(0,-1.5cm)$) {
            \textbf{Line 2:} Define the main method entry point
        };
    \end{tikzpicture}
    \end{center}
    \end{frame}

    \begin{frame}[fragile]{Code Execution Flow}
    \begin{center}
    \begin{tikzpicture}[remember picture, overlay]
        % Code listing (named node so we can reference its bounds)
        \node[anchor=center] (code) at (current page.center) {
            \hspace{1cm}
            \begin{minipage}{\textwidth}
    \begin{lstlisting}[style=JavaStyle,firstnumber=1]
    
public class Main {
    public static void main(String[] args) {
        System.out.println("Hello World: 1");
        System.out.println("Hello World: 2");
        System.out.println("Hello World: 3");
    \end{lstlisting}
            \end{minipage}
        };
        
        % === Spotlight the current line: draw a semi-opaque overlay with a hole ===
        \begin{scope}[on background layer]
            % Compute a window rectangle around the current line
            \path let \p1=(code.north west), \p2=(code.north east) in
                coordinate (winNW) at (\x1+0.9em, \y1-4.8em)  % top-left of the window
                coordinate (winSE) at (\x2-0.9em, \y1-3.7em); % bottom-right of the window
            
            % Wash out (pseudo-blur) everything except the window using even-odd rule
            \fill[black,opacity=0.5,even odd rule]
                (code.south west) rectangle (code.north east)
                (winNW) rectangle (winSE);
            
            % Draw the red box around the visible line
            \draw[red,very thick,rounded corners=2pt] (winNW) rectangle (winSE);
        \end{scope}
        
        % Explanation box below the code
        \node[anchor=north,fill=blue!10,draw=blue!50,thick,rounded corners,
            text width=0.8\textwidth,align=left,inner sep=10pt] 
            at ($(current page.north)+(0,-1.5cm)$) {
            \textbf{Line 3:} Print the first hello message
        };
    \end{tikzpicture}
    \end{center}
    \end{frame}

    \begin{frame}[fragile]{Code Execution Flow}
    \begin{center}
    \begin{tikzpicture}[remember picture, overlay]
        % Code listing (named node so we can reference its bounds)
        \node[anchor=center] (code) at (current page.center) {
            \hspace{1cm}
            \begin{minipage}{\textwidth}
    \begin{lstlisting}[style=JavaStyle,firstnumber=2]
    
    public static void main(String[] args) {
        System.out.println("Hello World: 1");
        System.out.println("Hello World: 2");
        System.out.println("Hello World: 3");
        System.out.println("Hello World: 4");
    \end{lstlisting}
            \end{minipage}
        };
        
        % === Spotlight the current line: draw a semi-opaque overlay with a hole ===
        \begin{scope}[on background layer]
            % Compute a window rectangle around the current line
            \path let \p1=(code.north west), \p2=(code.north east) in
                coordinate (winNW) at (\x1+0.9em, \y1-4.8em)  % top-left of the window
                coordinate (winSE) at (\x2-0.9em, \y1-3.7em); % bottom-right of the window
            
            % Wash out (pseudo-blur) everything except the window using even-odd rule
            \fill[black,opacity=0.5,even odd rule]
                (code.south west) rectangle (code.north east)
                (winNW) rectangle (winSE);
            
            % Draw the red box around the visible line
            \draw[red,very thick,rounded corners=2pt] (winNW) rectangle (winSE);
        \end{scope}
        
        % Explanation box below the code
        \node[anchor=north,fill=blue!10,draw=blue!50,thick,rounded corners,
            text width=0.8\textwidth,align=left,inner sep=10pt] 
            at ($(current page.north)+(0,-1.5cm)$) {
            \textbf{Line 4:} Print the second hello message
        };
    \end{tikzpicture}
    \end{center}
    \end{frame}

    \begin{frame}[fragile]{Code Execution Flow}
    \begin{center}
    \begin{tikzpicture}[remember picture, overlay]
        % Code listing (named node so we can reference its bounds)
        \node[anchor=center] (code) at (current page.center) {
            \hspace{1cm}
            \begin{minipage}{\textwidth}
    \begin{lstlisting}[style=JavaStyle,firstnumber=3]
    
        System.out.println("Hello World: 1");
        System.out.println("Hello World: 2");
        System.out.println("Hello World: 3");
        System.out.println("Hello World: 4");
        System.out.println("Hello World: 5");
    \end{lstlisting}
            \end{minipage}
        };
        
        % === Spotlight the current line: draw a semi-opaque overlay with a hole ===
        \begin{scope}[on background layer]
            % Compute a window rectangle around the current line
            \path let \p1=(code.north west), \p2=(code.north east) in
                coordinate (winNW) at (\x1+0.9em, \y1-4.8em)  % top-left of the window
                coordinate (winSE) at (\x2-0.9em, \y1-3.7em); % bottom-right of the window
            
            % Wash out (pseudo-blur) everything except the window using even-odd rule
            \fill[black,opacity=0.5,even odd rule]
                (code.south west) rectangle (code.north east)
                (winNW) rectangle (winSE);
            
            % Draw the red box around the visible line
            \draw[red,very thick,rounded corners=2pt] (winNW) rectangle (winSE);
        \end{scope}
        
        % Explanation box below the code
        \node[anchor=north,fill=blue!10,draw=blue!50,thick,rounded corners,
            text width=0.8\textwidth,align=left,inner sep=10pt] 
            at ($(current page.north)+(0,-1.5cm)$) {
            \textbf{Line 5:} Print the third hello message
        };
    \end{tikzpicture}
    \end{center}
    \end{frame}

    \begin{frame}[fragile]{Code Execution Flow}
    \begin{center}
    \begin{tikzpicture}[remember picture, overlay]
        % Code listing (named node so we can reference its bounds)
        \node[anchor=center] (code) at (current page.center) {
            \hspace{1cm}
            \begin{minipage}{\textwidth}
    \begin{lstlisting}[style=JavaStyle,firstnumber=4]
    
        System.out.println("Hello World: 2");
        System.out.println("Hello World: 3");
        System.out.println("Hello World: 4");
        System.out.println("Hello World: 5");
        System.out.println("Hello World: 6");
    \end{lstlisting}
            \end{minipage}
        };
        
        % === Spotlight the current line: draw a semi-opaque overlay with a hole ===
        \begin{scope}[on background layer]
            % Compute a window rectangle around the current line
            \path let \p1=(code.north west), \p2=(code.north east) in
                coordinate (winNW) at (\x1+0.9em, \y1-4.8em)  % top-left of the window
                coordinate (winSE) at (\x2-0.9em, \y1-3.7em); % bottom-right of the window
            
            % Wash out (pseudo-blur) everything except the window using even-odd rule
            \fill[black,opacity=0.5,even odd rule]
                (code.south west) rectangle (code.north east)
                (winNW) rectangle (winSE);
            
            % Draw the red box around the visible line
            \draw[red,very thick,rounded corners=2pt] (winNW) rectangle (winSE);
        \end{scope}
        
        % Explanation box below the code
        \node[anchor=north,fill=blue!10,draw=blue!50,thick,rounded corners,
            text width=0.8\textwidth,align=left,inner sep=10pt] 
            at ($(current page.north)+(0,-1.5cm)$) {
            \textbf{Line 6:} Print the fourth hello message
        };
    \end{tikzpicture}
    \end{center}
    \end{frame}

    \begin{frame}[fragile]{Code Execution Flow}
    \begin{center}
    \begin{tikzpicture}[remember picture, overlay]
        % Code listing (named node so we can reference its bounds)
        \node[anchor=center] (code) at (current page.center) {
            \hspace{1cm}
            \begin{minipage}{\textwidth}
    \begin{lstlisting}[style=JavaStyle,firstnumber=5]
    
        System.out.println("Hello World: 3");
        System.out.println("Hello World: 4");
        System.out.println("Hello World: 5");
        System.out.println("Hello World: 6");
        System.out.println("Hello World: 7");
    \end{lstlisting}
            \end{minipage}
        };
        
        % === Spotlight the current line: draw a semi-opaque overlay with a hole ===
        \begin{scope}[on background layer]
            % Compute a window rectangle around the current line
            \path let \p1=(code.north west), \p2=(code.north east) in
                coordinate (winNW) at (\x1+0.9em, \y1-4.8em)  % top-left of the window
                coordinate (winSE) at (\x2-0.9em, \y1-3.7em); % bottom-right of the window
            
            % Wash out (pseudo-blur) everything except the window using even-odd rule
            \fill[black,opacity=0.5,even odd rule]
                (code.south west) rectangle (code.north east)
                (winNW) rectangle (winSE);
            
            % Draw the red box around the visible line
            \draw[red,very thick,rounded corners=2pt] (winNW) rectangle (winSE);
        \end{scope}
        
        % Explanation box below the code
        \node[anchor=north,fill=blue!10,draw=blue!50,thick,rounded corners,
            text width=0.8\textwidth,align=left,inner sep=10pt] 
            at ($(current page.north)+(0,-1.5cm)$) {
            \textbf{Line 7:} Close the main method
        };
    \end{tikzpicture}
    \end{center}
    \end{frame}

    \begin{frame}[fragile]{Code Execution Flow}
    \begin{center}
    \begin{tikzpicture}[remember picture, overlay]
        % Code listing (named node so we can reference its bounds)
        \node[anchor=center] (code) at (current page.center) {
            \hspace{1cm}
            \begin{minipage}{\textwidth}
    \begin{lstlisting}[style=JavaStyle,firstnumber=6]
    
        System.out.println("Hello World: 4");
        System.out.println("Hello World: 5");
        System.out.println("Hello World: 6");
        System.out.println("Hello World: 7");
        System.out.println("Hello World: 8");
    \end{lstlisting}
            \end{minipage}
        };
        
        % === Spotlight the current line: draw a semi-opaque overlay with a hole ===
        \begin{scope}[on background layer]
            % Compute a window rectangle around the current line
            \path let \p1=(code.north west), \p2=(code.north east) in
                coordinate (winNW) at (\x1+0.9em, \y1-4.8em)  % top-left of the window
                coordinate (winSE) at (\x2-0.9em, \y1-3.7em); % bottom-right of the window
            
            % Wash out (pseudo-blur) everything except the window using even-odd rule
            \fill[black,opacity=0.5,even odd rule]
                (code.south west) rectangle (code.north east)
                (winNW) rectangle (winSE);
            
            % Draw the red box around the visible line
            \draw[red,very thick,rounded corners=2pt] (winNW) rectangle (winSE);
        \end{scope}
        
        % Explanation box below the code
        \node[anchor=north,fill=blue!10,draw=blue!50,thick,rounded corners,
            text width=0.8\textwidth,align=left,inner sep=10pt] 
            at ($(current page.north)+(0,-1.5cm)$) {
            \textbf{Line 8:} Close the Main class
        };
    \end{tikzpicture}
    \end{center}
    \end{frame}

    \begin{frame}[fragile]{Code Execution Flow}
    \begin{center}
    \begin{tikzpicture}[remember picture, overlay]
        % Code listing (named node so we can reference its bounds)
        \node[anchor=center] (code) at (current page.center) {
            \hspace{1cm}
            \begin{minipage}{\textwidth}
    \begin{lstlisting}[style=JavaStyle,firstnumber=7]
    
        System.out.println("Hello World: 5");
        System.out.println("Hello World: 6");
        System.out.println("Hello World: 7");
        System.out.println("Hello World: 8");
    }
    \end{lstlisting}
            \end{minipage}
        };
        
        % === Spotlight the current line: draw a semi-opaque overlay with a hole ===
        \begin{scope}[on background layer]
            % Compute a window rectangle around the current line
            \path let \p1=(code.north west), \p2=(code.north east) in
                coordinate (winNW) at (\x1+0.9em, \y1-4.8em)  % top-left of the window
                coordinate (winSE) at (\x2-0.9em, \y1-3.7em); % bottom-right of the window
            
            % Wash out (pseudo-blur) everything except the window using even-odd rule
            \fill[black,opacity=0.5,even odd rule]
                (code.south west) rectangle (code.north east)
                (winNW) rectangle (winSE);
            
            % Draw the red box around the visible line
            \draw[red,very thick,rounded corners=2pt] (winNW) rectangle (winSE);
        \end{scope}
        
        % Explanation box below the code
        \node[anchor=north,fill=blue!10,draw=blue!50,thick,rounded corners,
            text width=0.8\textwidth,align=left,inner sep=10pt] 
            at ($(current page.north)+(0,-1.5cm)$) {
            \textbf{Line 9:} Executing line 9
        };
    \end{tikzpicture}
    \end{center}
    \end{frame}

    \begin{frame}[fragile]{Code Execution Flow}
    \begin{center}
    \begin{tikzpicture}[remember picture, overlay]
        % Code listing (named node so we can reference its bounds)
        \node[anchor=center] (code) at (current page.center) {
            \hspace{1cm}
            \begin{minipage}{\textwidth}
    \begin{lstlisting}[style=JavaStyle,firstnumber=8]
    
        System.out.println("Hello World: 6");
        System.out.println("Hello World: 7");
        System.out.println("Hello World: 8");
    }
    \end{lstlisting}
            \end{minipage}
        };
        
        % === Spotlight the current line: draw a semi-opaque overlay with a hole ===
        \begin{scope}[on background layer]
            % Compute a window rectangle around the current line
            \path let \p1=(code.north west), \p2=(code.north east) in
                coordinate (winNW) at (\x1+0.9em, \y1-4.8em)  % top-left of the window
                coordinate (winSE) at (\x2-0.9em, \y1-3.7em); % bottom-right of the window
            
            % Wash out (pseudo-blur) everything except the window using even-odd rule
            \fill[black,opacity=0.5,even odd rule]
                (code.south west) rectangle (code.north east)
                (winNW) rectangle (winSE);
            
            % Draw the red box around the visible line
            \draw[red,very thick,rounded corners=2pt] (winNW) rectangle (winSE);
        \end{scope}
        
        % Explanation box below the code
        \node[anchor=north,fill=blue!10,draw=blue!50,thick,rounded corners,
            text width=0.8\textwidth,align=left,inner sep=10pt] 
            at ($(current page.north)+(0,-1.5cm)$) {
            \textbf{Line 10:} Executing line 10
        };
    \end{tikzpicture}
    \end{center}
    \end{frame}

    \begin{frame}[fragile]{Code Execution Flow}
    \begin{center}
    \begin{tikzpicture}[remember picture, overlay]
        % Code listing (named node so we can reference its bounds)
        \node[anchor=center] (code) at (current page.center) {
            \hspace{1cm}
            \begin{minipage}{\textwidth}
    \begin{lstlisting}[style=JavaStyle,firstnumber=9]
    
        System.out.println("Hello World: 7");
        System.out.println("Hello World: 8");
    }
    \end{lstlisting}
            \end{minipage}
        };
        
        % === Spotlight the current line: draw a semi-opaque overlay with a hole ===
        \begin{scope}[on background layer]
            % Compute a window rectangle around the current line
            \path let \p1=(code.north west), \p2=(code.north east) in
                coordinate (winNW) at (\x1+0.9em, \y1-4.8em)  % top-left of the window
                coordinate (winSE) at (\x2-0.9em, \y1-3.7em); % bottom-right of the window
            
            % Wash out (pseudo-blur) everything except the window using even-odd rule
            \fill[black,opacity=0.5,even odd rule]
                (code.south west) rectangle (code.north east)
                (winNW) rectangle (winSE);
            
            % Draw the red box around the visible line
            \draw[red,very thick,rounded corners=2pt] (winNW) rectangle (winSE);
        \end{scope}
        
        % Explanation box below the code
        \node[anchor=north,fill=blue!10,draw=blue!50,thick,rounded corners,
            text width=0.8\textwidth,align=left,inner sep=10pt] 
            at ($(current page.north)+(0,-1.5cm)$) {
            \textbf{Line 11:} Executing line 11
        };
    \end{tikzpicture}
    \end{center}
    \end{frame}

    \end{document}